% GENERATED FILE. Edit canonical lessons, then run npm run generate:books.
% canonical-source-hash: fnv1a64:3192a2616764d7ea
% canonical-lessons: ML-C103-pallikkoodam, ML-C103-paadam, ML-C103-pareeksha, ML-C103-maarkku, ML-R103-school-recall

\chapter{School}
\label{ch:ml-school}
\hlchaptermodality{103}

\begin{chapteropening}
\textbf{By the end of this chapter:} \emph{I can name a school and what happens in it — the lesson, the exam, and the marks.}

From cold: four school words, the three places they came from, and which one changes under a case ending.
\end{chapteropening}

\section[paḷḷikkūṭaṁ]{\ml{പള്ളിക്കൂടം} — a school}

\label{lesson:ML-C103-pallikkoodam}



\begin{quote}
Say \emph{a teacher}. Say \emph{a student}. You have had both of them for a long while, and never once had the building they are standing in.
\end{quote}

\subsection*{What to know first}
\textbf{\ml{പള്ളിക്കൂടം}} (\emph{paḷḷikkūṭaṁ}) — \textbf{a school}.

It is a long word, and it is long because it is \textbf{two words}. The first, \emph{paḷḷi}, is what Malayalam calls a \textbf{mosque or a church}. The second, \emph{kūṭaṁ}, is a \textbf{hall}. The ordinary everyday word for a school is a compound of two buildings.

\begin{culture}
That is not a coincidence, and it is not a metaphor either. \textbf{Teaching happened where the community already gathered} — the hall beside the place of worship was where children were taught, long before anybody built a separate building for it. The word kept the arrangement after the arrangement changed.

\textbf{Kerala has a second word too}, borrowed straight from English and heard constantly. This one is the native compound, and it is the one that tells you something.
\end{culture}

\subsection*{Your turn}
Say these aloud:

\begin{itemize}
\raggedright
  \item \emph{paḷḷikkūṭaṁ}, slowly — \emph{paḷḷi … kūṭaṁ}
  \item the two words it is made of, and what each one means
  \item the teacher, the student, and the building
\end{itemize}

\begin{tcolorbox}[breakable,colback=teal!4,colframe=teal!35!black,title={Before you move on}]
A school? (\textbf{\ml{പള്ളിക്കൂടം}}.) What are the two words inside it? (\emph{paḷḷi}, a place of worship, and \emph{kūṭaṁ}, a hall.) Who is in it that you could already name? (\textbf{\ml{അധ്യാപകൻ}} and \textbf{\ml{വിദ്യാർത്ഥി}}.)
\end{tcolorbox}
\section[pāṭhaṁ]{\ml{പാഠം} — a lesson}

\label{lesson:ML-C103-paadam}



\begin{quote}
Say \emph{a school}. Then say what happens to \textbf{-\ml{ം}} when a case ending is added to the word.
\end{quote}

\subsection*{What to know first}
\textbf{\ml{പാഠം}} (\emph{pāṭhaṁ}) — \textbf{a lesson}, the kind a school gives you: a chapter of a book, a thing to be learned.

The Sanskrit it comes from means \textbf{a recitation} — a saying-aloud. The lesson was named for the \textbf{saying}, not for the page it is written on, which is worth knowing about a word you are about to say aloud.

\begin{grammarlens}[title={Grammar Lens: you already know what it will do}]
\textbf{\ml{പാഠം} ends in -\ml{ം}}, and that tells you something before you are taught it. When a \textbf{case ending} is added, the \textbf{-\ml{ം}} gives way to \textbf{-\ml{ത്ത}-}:

\begin{quote}
\textbf{\ml{പാഠം}} $\to$ \textbf{\ml{പാഠത്തിൽ}} (\emph{pāṭhattil}) — \emph{in the lesson}
\end{quote}

\textbf{That is the same rule, unchanged, that the state name and the book followed.} You are not learning a new behaviour; you are meeting a new word that already belongs to a family you know.
\end{grammarlens}

\subsection*{Your turn}
Say these aloud:

\begin{itemize}
\raggedright
  \item \emph{pāṭhaṁ}
  \item \emph{pāṭhaṁ}, then \emph{pāṭhattil} — and name what changed
  \item what the Sanskrit behind it means
\end{itemize}

\begin{tcolorbox}[breakable,colback=teal!4,colframe=teal!35!black,title={Before you move on}]
A lesson? (\textbf{\ml{പാഠം}}.) What does its \textbf{-\ml{ം}} become when a case ending is added? (\textbf{-\ml{ത്ത}-} — \textbf{\ml{പാഠത്തിൽ}}.) What did the Sanskrit word behind it mean? (\textbf{A recitation}.)
\end{tcolorbox}
\section[parīkṣa]{\ml{പരീക്ഷ} — an exam}

\label{lesson:ML-C103-pareeksha}



\begin{quote}
Say \emph{a lesson}. Then say the apology you have had since early on — and listen for the \textbf{\ml{ക്ഷ}} in the middle of it.
\end{quote}

\subsection*{What to know first}
\textbf{\ml{പരീക്ഷ}} (\emph{parīkṣa}) — \textbf{an exam}.

\textbf{You can already write the hard part.} The \textbf{\ml{ക്ഷ}} in the middle is the one that opens that apology. It is not a new shape; it is a shape you have been saying since the beginning, sitting in a new word.

The Sanskrit behind it is \emph{pari-} (\textbf{around}) plus a root meaning \textbf{to look}. An examination is a \textbf{looking all around} something.

\begin{grammarlens}[title={Grammar Lens: this one does not end in -\ml{ം}}]
The word before this one ended in \textbf{-\ml{ം}} and told you what it would do. This one ends in \textbf{\ml{അ}}, and it does not join the same family:

\noindent
\begin{tabularx}{\linewidth}{@{}>{\raggedright\arraybackslash}X>{\raggedright\arraybackslash}X>{\raggedright\arraybackslash}X@{}}
\toprule
\textbf{} & \textbf{} & \textbf{} \\
\midrule
\textbf{\ml{പാഠം}} & \emph{pāṭhaṁ} & ends in \textbf{-\ml{ം}} \\
\textbf{\ml{പരീക്ഷ}} & \emph{parīkṣa} & ends in \textbf{\ml{അ}} \\
\bottomrule
\end{tabularx}

\textbf{Two words from the same chapter, in the same subject, behaving differently.} That is worth noticing rather than smoothing over: what a Malayalam word does when something attaches to it is decided by \textbf{how it ends}, not by what it means.
\end{grammarlens}

\subsection*{Your turn}
Say these aloud:

\begin{itemize}
\raggedright
  \item \emph{parīkṣa}
  \item the apology, then the exam — and name the shape they share
  \item \emph{pāṭhaṁ}, then \emph{parīkṣa}, and how each one ends
\end{itemize}

\begin{tcolorbox}[breakable,colback=teal!4,colframe=teal!35!black,title={Before you move on}]
An exam? (\textbf{\ml{പരീക്ഷ}}.) Which earlier word shares its middle shape? (That apology, \textbf{\ml{ക്ഷമിക്കണം}}.) Which of this chapter's two words ends in \textbf{-\ml{ം}}? (\textbf{\ml{പാഠം}} — not \textbf{\ml{പരീക്ഷ}}.)
\end{tcolorbox}
\section[mārkkŭ]{\ml{മാർക്ക്} — marks}

\label{lesson:ML-C103-maarkku}



\begin{quote}
Say \emph{an exam}. Then count in tens — you will want the numbers in a moment.
\end{quote}

\subsection*{What to know first}
\textbf{\ml{മാർക്ക്}} (\emph{mārkkŭ}) — \textbf{marks}, what an exam gives back to you.

\textbf{Say it out loud and you will hear English.} It is the English word \emph{mark}, borrowed whole. Kerala's schools took the English examination system and the English word along with it, and this is the word people reach for.

\begin{grammarlens}[title={Grammar Lens: what the borrowing had to change}]
The word is English, and the \textbf{spelling is not}. Two things happened to it on the way in, and you can see both:

\begin{itemize}
\raggedright
  \item the \textbf{\ml{ർ}} before a doubled \textbf{\ml{ക്ക}} — the shape you already write in \textbf{\ml{ശർക്കര}} and \textbf{\ml{ചേർക്കുക}};
  \item the \textbf{half-u} at the end, the faint vowel that closes so many Malayalam words and that English \emph{mark} does not have at all.
\end{itemize}

\textbf{A borrowed word does not stay foreign.} It is made to end the way Malayalam words end, and then it behaves like one.
\end{grammarlens}

\subsection*{Your turn}
Say these aloud:

\begin{itemize}
\raggedright
  \item \emph{mārkkŭ}, and hear the English inside it
  \item \emph{parīkṣa}, then \emph{mārkkŭ} — the exam, then what it gives back
  \item a number in tens, then \emph{mārkkŭ}
\end{itemize}

\begin{tcolorbox}[breakable,colback=teal!4,colframe=teal!35!black,title={Before you move on}]
Marks? (\textbf{\ml{മാർക്ക്}}.) Where does the word come from? (\textbf{English} — \emph{mark}.) What two things did Malayalam do to its spelling? (\textbf{The chillu \ml{ർ} before the doubled \ml{ക്ക}}, and \textbf{the half-u at the end}.)
\end{tcolorbox}
\section[paḷḷikkūṭaṁ enna ōrmma]{(school, lesson, exam, marks) — from cold}

\label{lesson:ML-R103-school-recall}



\begin{quote}
Close the lessons before this one. Say \emph{school}, \emph{lesson}, \emph{exam} and \emph{marks}.
\end{quote}

\begin{grammarlens}[title={Grammar Lens: four words, three different origins}]
One small corner of one subject, and the words in it came from three places:

\noindent
\begin{tabularx}{\linewidth}{@{}>{\raggedright\arraybackslash}X>{\raggedright\arraybackslash}X>{\raggedright\arraybackslash}X@{}}
\toprule
\textbf{} & \textbf{} & \textbf{} \\
\midrule
\textbf{\ml{പള്ളിക്കൂടം}} & \emph{paḷḷikkūṭaṁ} & a compound of two buildings \\
\textbf{\ml{പാഠം}}, \textbf{\ml{പരീക്ഷ}} & \emph{pāṭhaṁ}, \emph{parīkṣa} & Sanskrit \\
\textbf{\ml{മാർക്ക്}} & \emph{mārkkŭ} & English \\
\bottomrule
\end{tabularx}

\textbf{That is what a school vocabulary looks like in Malayalam}, and it is not untidy — it is a record. The building is named the old way, the scholarly words came with the scholarly tradition, and the examination system arrived with its own word attached.
\end{grammarlens}

\begin{grammarlens}[title={What you've built}]
\textbf{You had the people before you had the place.} The teacher and the student have been yours for a long while with nothing around them; now there is a building, a thing taught in it, a test of the thing, and a number at the end.

And the two Sanskrit words do \textbf{not} behave alike: \textbf{\ml{പാഠം}} ends in \textbf{-\ml{ം}} and gives way to \textbf{-\ml{ത്ത}-} when a case ending arrives; \textbf{\ml{പരീക്ഷ}} ends in \textbf{\ml{അ}} and does no such thing. Same subject, same origin, different endings — and the ending is what decides.
\end{grammarlens}

\subsection*{Your turn}
Say these aloud:

\begin{itemize}
\raggedright
  \item \emph{paḷḷikkūṭaṁ}, \emph{pāṭhaṁ}, \emph{parīkṣa}, \emph{mārkkŭ}
  \item which of the four is English, and which two are Sanskrit
  \item \emph{pāṭhaṁ}, then \emph{pāṭhattil}
  \item the teacher and the student, now inside the building
\end{itemize}

\begin{tcolorbox}[breakable,colback=teal!4,colframe=teal!35!black,title={Before you move on}]
School, exam, marks? (\textbf{\ml{പള്ളിക്കൂടം}, \ml{പരീക്ഷ}, \ml{മാർക്ക്}}.) Which one is English? (\textbf{\ml{മാർക്ക്}}.) Which word in the chapter changes when a case ending arrives, and to what? (\textbf{\ml{പാഠം}} $\to$ \textbf{\ml{പാഠത്തിൽ}}.)
\end{tcolorbox}
